\documentclass[11pt]{article}

\usepackage[a4paper,margin=30mm]{geometry}
\usepackage{amsmath,amssymb,amsthm,mathtools}
\usepackage[hidelinks]{hyperref}

\newtheorem{theorem}{Theorem}[section]
\newtheorem{proposition}[theorem]{Proposition}
\newtheorem{lemma}[theorem]{Lemma}
\newtheorem{corollary}[theorem]{Corollary}
\theoremstyle{definition}
\newtheorem{definition}[theorem]{Definition}
\theoremstyle{remark}
\newtheorem{remark}[theorem]{Remark}

\newcommand{\N}{\mathbf N}
\newcommand{\Z}{\mathbf Z}
\newcommand{\cS}{\mathcal S}
\newcommand{\cG}{\mathcal G}

\title{A Machine-Certified Closure of Erd\H{o}s Problem \#1212}
\author{}
\date{Public minimal proof, 14 July 2026}

\begin{document}
\maketitle

\begin{abstract}
We prove the original statement of Erd\H{o}s Problem \#1212: the positive
visible lattice, after deletion of every prime--prime vertex, contains an
unbounded unit-step path.  The proof first reduces the path to tightness of
bad separating contours.  Fixed-contour prime--prime contributions have zero
translation density, so compatible prime-divisibility witnesses give a
positive, coefficientwise necklace expansion.  Splitting that expansion at
the first repeated prime gives an all-fresh term and a repeat operator.  The
all-fresh term is summable.  A certified tilted Hashimoto inequality bounds
all repeated primes at least $61$ by $0.067$.  For the twelve remaining
primes, exact Fourier projection, exact-rational root enclosures, binary64
centres, directed-rounding tails, and a full-state interval comparison prove
the common-cone bound
\[
 \mathcal R_{<61}u\leq U
 \leq0.91717912709094451u<0.933u.
\]
The comparison covers all $55{,}076{,}704$ allocated states and has zero
threshold, positive-over-zero, negative-upper, and nonfinite failures.  Thus
the complete first-repeat norm is strictly below one, separating contour mass
cannot escape, and the required path exists.  This note contains only that
closed proof chain and its machine-certificate contract.
\end{abstract}

\section{The unchanged target}

Let
\[
 \cG=\{(x,y)\in\N^2:\gcd(x,y)=1\}
\]
with an edge when exactly one coordinate changes by $1$ or $-1$, and put
\[
 \cS=\{(x,y)\in\cG:x>1,\ y>1,
       \text{ and at least one of $x,y$ is composite}\}.
\]
Thus $\cS$ is obtained from the positive visible lattice by deleting the two
coordinate boundary strips and every prime--prime vertex.

\begin{theorem}[Erd\H{o}s \#1212]
\label{thm:main-target}
There is a sequence $(P_t)_{t\geq0}$ in $\cS$ such that consecutive terms are
adjacent and, for every $R\geq1$, some $P_t$ has a coordinate at least $R$.
\end{theorem}

\begin{proposition}[Finite reachability dictionary]
\label{prop:finite-reachability}
Theorem~\ref{thm:main-target} is equivalent to the existence of one
$v\in\cS$ from which finite paths in $\cS$ leave every finite box.
\end{proposition}

\begin{proof}
An infinite path gives the finite paths by truncation.  Conversely, remove
loops from each finite path.  The rooted tree of finite simple paths beginning
at $v$ is locally finite and has arbitrarily large depth.  K\"onig's infinity
lemma gives an infinite ray, which cannot remain in a finite box.
\end{proof}

\section{Fixed contours and the only escape event}

Put
\[
 \mathcal D=\{(x,y)\in\N^2:\gcd(x,y)>1\},\qquad
 \mathcal P^{(2)}=\{(p,q):p,q\text{ are prime}\}.
\]
Away from the boundary strips, $\N^2\setminus\cS=\mathcal D\cup
\mathcal P^{(2)}$.

Let $A\subset\Z^2$ be finite.  A map from $A$ to the primes is
\emph{compatible} when equal labels $p$ occur only at sites congruent
coordinatewise modulo $p$.  Give such a map the weight
\[
 \omega(w)=\prod_{p\in w(A)}p^{-2},\qquad
 Z(A)=\sum_{w\ \mathrm{compatible}}\omega(w).
\]
The sum is finite: a fibre of size at least two can be labelled only by a
prime dividing both coordinate differences of two sites, while singleton
fibres contribute the convergent sum $\sum_p p^{-2}$.

\begin{proposition}[Exact witness expansion]
\label{prop:fixed-witness-expansion}
For
\[
 H_A(T)=\{z\in[1,T]^2:z+a\in\mathcal D\text{ for every }a\in A\},
\]
one has
\[
 \limsup_{T\to\infty}\frac{|H_A(T)|}{T^2}\leq Z(A).
\]
A fixed compatible witness $w$ is realized with natural density exactly
$\omega(w)$.
\end{proposition}

\begin{proof}
Choose one common prime divisor at every translated site.  Equal chosen primes
force compatibility.  Conversely, compatibility and the Chinese remainder
theorem give one residue class in each coordinate modulo
$\prod_{p\in w(A)}p$, of density $\omega(w)$.  Restricting first to primes at
most $R$, applying the union bound, and then sending $R$ to infinity justifies
the infinite prime sum.
\end{proof}

\begin{theorem}[Fixed-contour prime cancellation]
\label{thm:fixed-contour-cancellation}
Let $B_A(T)$ count translations $z\in[1,T]^2$ for which every $z+a$ is
outside $\cS$.  Boundary translations contribute only $O_A(T)$, and
\[
 \frac{|B_A(T)|-|H_A(T)|}{T^2}\longrightarrow0.
\]
Thus the translation density of a fixed bad contour is exactly the density of
the same contour made entirely of nonvisible sites.
\end{theorem}

\begin{proof}
The inclusion $H_A(T)\subseteq B_A(T)$ is immediate.  A translation in the
difference has at least one prime--prime site.  For fixed $a$ there are
$O_A(T^2/\log^2T)$ such translations, and a union bound over $A$ proves the
claim.
\end{proof}

Let $C_\infty$ be Vardi's unique infinite component of the full visible
lattice; it has a positive asymptotic density $\theta$ \cite{Vardi1999}.  Let
$\mathfrak C$ be the simple closed star-contours surrounding the origin.  For
$M\geq1$ let $E_{T,M}$ be the set of
$z\in C_\infty\cap[1,T]^2\cap\cS$ for which some
$\gamma\in\mathfrak C$ with $|\gamma|>M$ satisfies
$z+\gamma\subseteq\N^2\setminus\cS$, and put
\[
 \tau(M)=\limsup_{T\to\infty}\frac{|E_{T,M}|}{T^2}.
\]

\begin{theorem}[No escape closes the problem]
\label{thm:no-escape-closes}
If $\tau(M)\to0$ as $M\to\infty$, then Theorem~\ref{thm:main-target} holds.
\end{theorem}

\begin{proof}
Assume every component of $\cS$ is finite.  The planar outer-boundary lemma
gives a simple closed bad star-contour around every $z\in\cS$.  For fixed
$M$, there are finitely many relative contour shapes of length at most $M$.
A short contour around a point of $C_\infty$ cannot consist entirely of sites
in $\mathcal D$, since an infinite visible path from that point must cross it.
It therefore contains a prime--prime site.  By
Theorem~\ref{thm:fixed-contour-cancellation}, the roots admitting any of the
finitely many short shapes have density zero.  Prime--prime roots also have
density zero, while $C_\infty\cap\cS$ has density $\theta>0$.  Hence
$\theta\leq\tau(M)$ for every $M$, a contradiction.  Thus $\cS$ has an
unbounded component, and Proposition~\ref{prop:finite-reachability} extracts
the path.
\end{proof}

\section{Positive first-repeat renewal}

Fix $Y=11$, put $Q=2310$, and let
\[
 F_{11}=\{(a,b):\gcd(a,b,Q)>1\}.
\]
Contract only the finite star-components of $F_{11}$ and retain every site
outside $F_{11}$, together with all component--site incidences and external
star-edges.  On directed incidences use the Hashimoto update that forbids the
immediate reverse edge and charges a factor $q$ upon entering an external
site.  Write
\[
 \mathsf H(q)=\mathsf A+q\mathsf B,
\]
where $\mathsf B$ consists of transitions entering an external site.

For distinct tail primes define
\[
 d_m=\sum_{\substack{p_1,\ldots,p_m>11\\p_i\ \mathrm{distinct}}}
       \frac1{p_1^2\cdots p_m^2}.
\]

\begin{lemma}[All-fresh summability]
\label{lem:fresh-prime-summability}
For every $C>0$, $\sum_{m\geq0}C^m d_m<\infty$.
\end{lemma}

\begin{proof}
Let $e_m=d_m/m!$.  For $t>0$,
\[
 \sum_{m\geq0}e_mt^m
 \leq\prod_{n\geq1}\left(1+\frac{t}{n^2}\right)
 =\frac{\sinh(\pi\sqrt t)}{\pi\sqrt t}
 \leq e^{\pi\sqrt t}.
\]
Thus $e_m\leq e^{\pi\sqrt t}t^{-m}$.  Taking
$\sqrt t=2m/\pi$ and using $m!\leq m^m$ gives
\[
 d_m\leq\left(\frac{e^2\pi^2}{4m}\right)^m,
\]
which is superexponentially summable against $C^m$.
\end{proof}

Read the tail-prime labels in order and stop at the first repeated label.
Let $\mathcal F$ be the all-fresh Green function and $\mathcal R$ the positive
operator of compatible first-repeat excursions, retaining the lifted
displacement.  The complete necklace Green function satisfies the
coefficientwise renewal inequality
\begin{equation}
\label{eq:first-repeat-renewal}
 \mathcal G\preccurlyeq\mathcal F+\mathcal R\mathcal G.
\end{equation}
No cancellation is used.

\begin{lemma}[Subcritical renewal implies no escape]
\label{lem:renewal-no-escape}
If $\mathcal F$ is finite in the rooted area-weighted necklace norm and
$\|\mathcal R\|<1$ in the same positive norm, then $\tau(M)\to0$.
\end{lemma}

\begin{proof}
Iterating \eqref{eq:first-repeat-renewal} coefficientwise gives
\[
 \mathcal G\preccurlyeq\sum_{j\geq0}\mathcal R^j\mathcal F,
 \qquad
 \|\mathcal G\|\leq\frac{\|\mathcal F\|}{1-\|\mathcal R\|}<\infty.
\]
Every compatible translated bad contour contributes to the nonnegative
witness expansion bounded by this necklace series.  The contribution of
contours of length greater than $M$ is the tail of a finite nonnegative
series, hence tends to zero.  Proposition~\ref{prop:fixed-witness-expansion}
turns the same bound into translation upper density, while
Theorem~\ref{thm:fixed-contour-cancellation} removes fixed-shape
prime--prime terms.  Therefore the defining upper density $\tau(M)$ tends to
zero.
\end{proof}

\section{The infinite repeat-prime tail}

Put
\[
 q_{\rm tr}=\sum_{p>11}\frac1{p^2},\qquad
 \bar q=\frac{22751}{10^6}.
\]

\begin{lemma}[Translation-tail majorant]
\label{lem:translation-tail-majorant}
One has $q_{\rm tr}<\bar q$.
\end{lemma}

\begin{proof}
Every prime greater than $11$ is coprime to $2310$, and the only integer in
$[1,11]$ coprime to $2310$ is $1$.  Therefore
\[
 q_{\rm tr}\leq
 \frac{\pi^2}{6}\prod_{p\leq11}\left(1-\frac1{p^2}\right)-1
 \leq\frac{1722359}{75707401}<\frac{22751}{10^6},
\]
where the middle inequality uses $\pi<355/113$.
\end{proof}

For a lifted displacement $d$ and $\xi\in\mathbb R^2$, let
$\mathsf H_\xi(q)$ multiply the transition by $e^{\langle\xi,d\rangle}$.

\begin{theorem}[Certified translation kernel]
\label{thm:certified-translation-kernel}
On the cyclic support of the $2310$-period incidence graph,
\[
 \rho(\mathsf H(\bar q))<0.97,
 \qquad
 \rho(\mathsf H_\xi(\bar q))<0.97
\]
for $\xi=(\pm3/20,0),(0,\pm3/20)$.
\end{theorem}

\begin{proof}
The finite graph has $333{,}280$ permanent-component types,
$3{,}317{,}760$ external sites, $12{,}387{,}904$ directed
component--site states, and $14{,}407{,}568$ directed external-site states.
The largest permanent component has $77$ sites and coordinate span $12$.

Power iteration only constructs a positive witness $v$.  Acceptance is the
statewise Collatz inequality
\[
 \mathsf H_\xi(\bar q)v\leq0.97v.
\]
Every binary64 source sum is recomputed in long-double arithmetic, enlarged
by a relative $10^{-10}$, and compared with a target reduced by the same
amount.  Tilt factors are rational upper enclosures obtained from Taylor
intervals.  The untilted maximum upper ratio is $0.963345$ and the
$(3/20,0)$ ratio is $0.965210$, with zero failures.  Reflections and coordinate
interchange give the remaining tilts.  The source hash and exact transcript
are frozen in \path{certificate/hashimoto.txt}; the checker is
\path{kernel/skeleton_winding_kernel.cpp}.
\end{proof}

Let $\mathcal R_{\geq61}$ be the first-repeat contribution whose repeated
label is a prime at least $61$.

\begin{theorem}[Large repeat primes]
\label{thm:large-repeat-tail}
In the certificate norm of
Theorem~\ref{thm:certified-translation-kernel},
\[
 \|\mathcal R_{\geq61}\|<0.067.
\]
\end{theorem}

\begin{proof}
Two sites carrying the same prime $p$ differ by a nonzero vector in
$p\Z^2$.  One of the four oriented tilts contributes at most $e^{-3p/20}$.
Removing both scalar endpoint charges conservatively gives
\[
 \|\mathcal R_{\geq61}\|
 \leq\frac4{\bar q^2(1-0.97)}
       \sum_{p\geq61}\frac{e^{-3p/20}}{p^2}.
\]
Replacing primes by all integers and summing geometrically,
\[
 \sum_{p\geq61}\frac{e^{-3p/20}}{p^2}
 \leq\frac{e^{-183/20}}{61^2(1-e^{-3/20})}
 <\frac{23}{24000\cdot61^2}.
\]
Consequently
\[
 \|\mathcal R_{\geq61}\|
 <\frac{1150000000000}{17334174345489}<0.067.
\]
\end{proof}

\section{The finite twelve-prime certificate}

Only
\[
 \{13,17,19,23,29,31,37,41,43,47,53,59\}
\]
remain.  Decompose them as
\[
 S_0=\{13,17,19,23,29\},\qquad
 S_1=\{31,37,41,43,47,53,59\},\qquad
 s_i=\sum_{p\in S_i}p^{-2}.
\]

For a character $\chi$ of $(\Z/p\Z)^2$, twist every lifted displacement and
write the operators as $\mathsf H_\chi(\bar q)$ and $\mathsf B_\chi$.
Character orthogonality gives the exact projector
\[
 \frac1{p^2}\sum_\chi\chi(d)=\mathbf1_{d\in p\Z^2}.
\]
Together with the common-prime cost, the coefficient is $p^{-4}$.

Let $\eta_m=d_m/\bar q^m$.  It is decreasing, because
$d_{m+1}\leq\bar qd_m$.  For nonnegative coefficients $c_n$ define
\[
 \mathcal Q_{S,N}(c)=
 \sum_{p\in S}\frac1{p^4}\sum_\chi
 \mathsf B_\chi\left(\sum_{n=0}^Nc_n
 \mathsf H_\chi(\bar q)^n\right)\mathsf B,
\]
where $c_n=1$ for $n\leq20$ and
$c_n=\eta_{\lfloor n/2\rfloor}$ thereafter.

The common-cone vector $u$ is strictly positive on every used state.  Let
$C_i$ be the binary64 centre of $\mathcal Q_{S_i,N_i}(c)u$, with
$N_0=40$ and $N_1=24$, and define the directed-rounding positive images
\begin{align*}
 A_i&=\mathsf B\left(\sum_{n=0}^{N_i}c_n
          \mathsf H(\bar q)^n\right)\mathsf Bu,\\
 T_{25}&=\mathsf B\left(\sum_{n=25}^{400}
          \eta_{\lfloor n/2\rfloor}\mathsf H(\bar q)^n\right)\mathsf Bu,\\
 T_{41}&=\mathsf B\left(\sum_{n=41}^{400}
          \mathsf H(\bar q)^n\right)\mathsf Bu.
\end{align*}

All $16{,}548$ characters are included.  The root tables are verified without
floating trigonometry: Machin's identity and rational alternating arctangent
series enclose $\pi$, followed by quadrant reduction and rational Taylor
intervals for sine and cosine.  With $\epsilon=2^{-24}$, every stored root
component is within $4\epsilon$, hence every root has Euclidean error below
$6\epsilon$.  A retained word has at most $41$ phase factors and
\[
 \frac{(1+6\epsilon)^{41}-1}{\epsilon}<247.
\]
Fewer than $20{,}000$ binary64 rounding units occur along a state path and
complete projection, while $\gamma_{20000}(2^{-53})<\epsilon$.  Therefore
\begin{equation}
\label{eq:phase-error-enclosure}
 |\mathcal Q_{S_i,N_i}(c)u-C_i|\leq500\epsilon s_iA_i.
\end{equation}

The post-$400$ transfer remains in the same cone.  With the certified
untilted witness $V$, directed comparison gives
\[
 \mathsf Bu\leq3067V,\qquad \mathsf BV\leq14400u,
 \qquad \mathsf H(\bar q)V\leq0.97V.
\]
Since $s_0+s_1<0.01923$ and $(0.97)^{401}<5\cdot10^{-6}$, the post-$400$
contribution is below $1.14\cdot10^{-6}u$.  Define
\begin{equation}
\label{eq:repeat-directed-upper}
 U=C_0+C_1+500\epsilon(s_0A_0+s_1A_1)
   +s_1T_{25}+8\cdot10^{-9}s_0T_{41}
   +1.14\cdot10^{-6}u.
\end{equation}

\begin{theorem}[Twelve-prime directed interval certificate]
\label{thm:repeat-twelve-interval}
For the frozen vector $u$,
\[
 \mathcal R_{<61}u\leq U<0.933u.
\]
More precisely,
\[
 \max_{u(x)>0}\frac{U(x)}{u(x)}
 \leq0.91717912709094451<0.933,
\]
and $U(x)=0$ at every allocated coordinate with $u(x)=0$.
\end{theorem}

\begin{proof}
The Fourier centres are enclosed by
\eqref{eq:phase-error-enclosure}; $T_{25}$ and $T_{41}$ are directed upper
images; the last term of \eqref{eq:repeat-directed-upper} covers the infinite
tail.  The final CUDA kernel uses upward binary64 operations for every
positive sum, product and division, and compares with the downward binary64
neighbour of $0.933$.  It checks all $55{,}076{,}704$ allocated coordinates.
The threshold-failure, positive-over-zero, negative-upper, and nonfinite
counts are all zero.  The maximum is the displayed value.  Every source,
input, output, header, and length is bound by
\path{certificate/manifest.json}; the frozen output is
\path{certificate/admission.txt}.
\end{proof}

\section{Assembly}

\begin{theorem}[Closure of Erd\H{o}s Problem \#1212]
\label{thm:erdos1212-close}
Theorem~\ref{thm:main-target} holds.
\end{theorem}

\begin{proof}
Lemma~\ref{lem:fresh-prime-summability} and
Theorem~\ref{thm:certified-translation-kernel} make the all-fresh term finite
in the positive necklace norm.  The repeat operator is the disjoint sum of
its $p<61$ and $p\geq61$ parts.  Theorems
\ref{thm:repeat-twelve-interval} and \ref{thm:large-repeat-tail} give
\[
 \|\mathcal R\|<0.933+0.067=1.
\]
Lemma~\ref{lem:renewal-no-escape} now proves $\tau(M)\to0$ without moving the
first-repeat cut.  Theorem~\ref{thm:no-escape-closes} and
Proposition~\ref{prop:finite-reachability} yield the unbounded path in $\cS$.
\end{proof}

\section{Certificate contract}

The certificate kernel consists of one periodic-graph/Collatz checker, two
directed image generators, the final directed interval checker, and the
exact-rational root verifier.  The two root tables are committed because they
are small.  The dense state images are excluded from Git because each double
image is $440{,}613{,}656$ bytes; their exact byte lengths and SHA-256 digests
are frozen in \path{certificate/manifest.json}.  The repository gate rejects
unlisted files, missing theorem bindings, altered kernel sources, altered root
tables, and any admission transcript without all four zero-failure counts.
With the hash-pinned dense payload present, the full replay script rebuilds
the final CUDA checker and reruns the statewise comparison.

\begin{thebibliography}{3}

\bibitem{Erdos1980}
P.~Erd\H{o}s,
\emph{A survey of problems in combinatorial number theory},
Annals of Discrete Mathematics 6 (1980), 89--115.

\bibitem{Vardi1999}
I.~Vardi,
\emph{Deterministic percolation},
Communications in Mathematical Physics 207 (1999), 43--66.

\bibitem{Bloom1212}
T.~F. Bloom,
\emph{Erd\H{o}s Problem \#1212},
\url{https://www.erdosproblems.com/1212}.

\end{thebibliography}

\end{document}
